\documentclass[12pt]{article}
\usepackage[french]{babel}
\usepackage{url}
\usepackage[T1]{fontenc}
\usepackage{graphicx}
\graphicspath{{images/}}
\usepackage{parskip}
\usepackage{fancyhdr}
\usepackage{vmargin}
\usepackage{xcolor}
\usepackage{bbm}
\usepackage{amsmath,amssymb}
\usepackage{amsthm}
\usepackage{dsfont}
\usepackage{stmaryrd}
\usepackage{systeme}
\usepackage{enumitem}
\usepackage{xcolor}
\usepackage{pifont}
\usepackage[cache=false]{minted}
\definecolor{LightGray}{gray}{0.95}

\title{Compte-rendu}
\author{PIARD A. - JACQUET R.}
\date{\today}

\makeatletter
\let\thetitle\@title
\let\theauthor\@author
\let\thedate\@date
\makeatother

\pagestyle{fancy}
\fancyhf{}
\rhead{\theauthor}
\lhead{\thetitle}
\cfoot{\thepage}
\def\dotfill#1{\cleaders\hbox to #1{.}\hfill}
\newcommand\dotline[2][.5em]{\leavevmode\hbox to #2{\dotfill{#1}\hfil}}

\begin{document}

\begin{titlepage}
	\centering
    \vspace*{0.5 cm}
    \textsc{\LARGE Droit et Conduite de Projet . 2020-2021}\\[1.0 cm]
    \dotline[15pt]{15cm}\\
	\includegraphics[scale = 2.2]{logo.png}
	\dotline[15pt]{15cm}\\
	\vspace{1.5cm}
	\textsc{\Large Faculté des Sciences et Techniques}\\
	\textsc{\large Master 1 - Maths. CRYPTIS}\\[1.0 cm]
	\rule{\linewidth}{0.2 mm} \\[0.4 cm]
	{ \huge \bfseries \color{blue} \thetitle}\\
	\rule{\linewidth}{0.2 mm} \\[1.5 cm]
	
	\begin{minipage}{0.4\textwidth}
		\begin{flushleft} \large
			\emph{A l'attention de :}\\
			M. CRESPIN\\
			M. DUSART\\
			M. CONCHON\\
		\end{flushleft}
	\end{minipage}
	\begin{minipage}{0.5\textwidth}
    	\begin{flushright} \large
		\emph{Rédigé par :}\\
		PIARD A.\\
		JACQUET R.\\
		\phantom{a}\\
		\end{flushright}
	\end{minipage}\\[2 cm]
\end{titlepage}

\tableofcontents
\pagebreak

\section{Introduction}
Le but de notre projet était d'entraîner et de tester les compétences, en terme de programmation, de nos collègues du \textsf{Master 1 Maths CRYPTIS}. Comme énoncé dans notre fiche d'avant-projet, deux \textsl{Competitive Programming Test} étaient prévus, avec les quatre mêmes participants, et nous avions deux protocoles organisationnelles qui s'offraient à nous en vue de la condition sanitaire actuelle. Le protocole ayant été retenu est celui du distanciel.\\
\\
Le premier sujet a été traité le 1er février 2021 et le second le 8 février 2021.

\vfill \eject

\section{Premières réactions}
À l'issue du premier \textsl{Competitive Programming Test}, nous nous sommes rendus compte que les étudiants ont mis beaucoup plus de temps que prévu pour faire les exercices, ce qui fait qu'aucun d'eux n'a terminé le premier sujet. Il se trouve que les exercices n'étaient pas triés par ordre de difficulté, mais que les étudiants ont tous choisis de commencer par le premier exercice.\\\\
À l'issue du premier \textsl{Competitive Programming Test} et en vue du second  qui était prévu la semaine suivante, nous avons interrogés les participants pour les aider à comprendre ce qui n'allait pas. Celui-ci s'est bien mieux passé. Les étudiants n'ont pas terminés le sujet mais ils ont tous traités au moins un exercice. La plupart des étudiants s'étant intéressés à deux exercices.

\section{Analyse des sujets}
\subsection{Premier sujet}
\subsubsection{Premier exercice}
%Joindre image du sujet
%Dire ce qui semble les avoir gênés, nos réponses apportées, etc

%--- Objectifs du sujet
%--- Comment il a été traité dans l'ensemble
%--- Conclusion de l'expérience de ce sujet

L'objectif principal de cet exercice était de mettre l'étudiant face à un chiffré dont la méthode de chiffrement est inconnue. Il a été utilisé deux clés pour chiffrer le message. L'exercice en lui-même n'était pas difficile mais il fallait que l'étudiant tente diverses choses pour comprendre comment le message avait été chiffré. Le chiffrement en question est un chiffrement par substitution (de type César) pour lequel les clés étaient \textsf{3} et \textsf{13}. Malgré les indications la majorité des étudiants a perdu beaucoup de temps sur cet exercice, ce qui a influé sur l'ensemble du sujet. En effet un étudiant n'a pas pu effectuer l'exercice 2 et n'a pas trouvé comment résoudre l'exercice 1.


\subsubsection{Second exercice}

L'objectif de ce second exercice était de développer la méthode de déchiffrement d'un chiffrement précisément décrit. On rappelle que le choix du langage de programmation n'était pas imposé et que suivant le langage utilisé cela est plus ou moins rapide à écrire. Les étapes de l'algorithme étaient bien décomposées mais malgré tout nous n'avons pas eu beaucoup de retours, en raison du temps mal géré par les étudiants lors de l'exercice 1.

\pagebreak

\subsection{Second sujet}
\subsubsection{Premier exercice}
L'objectif de cet exercice était de programmer l'algorithme de chiffrement écrit en français dans l'énoncé. Les seules difficultés étaient la maîtrise des opérateurs logiques et la conversion de type. Contrairement à nos attentes, cet exercice n'a pas très bien été réussi. Étant donné que l'épreuve a été conduite en distanciel, les étudiants avaient accès à toutes les ressources possibles. Malgré tout, certains n'ont pas su trouver comment effectuer un \textsf{ET Logique}.

\subsubsection{Second exercice}
L'objectif de cet exercice était d'écrire la méthode de déchiffrement à partir de celle qui a permit le chiffrement. Les opérations étaient donc à inverser. Si les étudiants possédaient la bonne fonction, ils leur étaient alors possible de déchiffrer le message. Ici aussi, de nombreuses difficultés se sont présentées aux étudiants, comme le passage de tableaux de bits à une chaîne de caractères.

\subsubsection{Troisième exercice}
L'objectif de cet exercice était de déchiffrer un fichier binaire donné. Afin de déchiffrer ce fichier il fallait trouver la clef de chiffrement, qui était écrite sur l'image dont l'extension n'est pas la bonne. Il fallait donc au préalable trouver le bon format d'image avec la commande type \textsf{xxd}, ou \textsf{exiftool}, ou autre. Une fois la clef récupérée, il fallait la concaténer autant de fois que nécessaire pour obtenir une taille de clef similaire à la taille du fichier binaire. Enfin, il restait simplement à effectuer un \textsf{XOR} entre le fichier binaire et la clef pour obtenir l'image finale. Cet exercice était volontairement légèrement plus difficile et comme attendu il n'a pas été traité.

\pagebreak

\section{Comparaisons entre étudiants}
\subsection{Sujet n°1}
\subsubsection{Premier exercice}

\begin{minted}
[
frame=lines,
framesep=2mm,
baselinestretch=1.2,
bgcolor=LightGray,
fontsize=\footnotesize,
linenos
]
{python}
import re

chiffre = "IPLU'QV\Ð 'Ç'[VP3'SL'TLZZHNL'JHJOÐ'LZ['A'HGT1GXJ'("

n = len(chiffre)
print(n)
\end{minted}
Ici il est clair que l'étudiant n'a eu aucune idée de comment traiter l'exercice. Du moins il n'en a pas laissé de traces.

\begin{minted}
[
frame=lines,
framesep=2mm,
baselinestretch=1.2,
bgcolor=LightGray,
fontsize=\footnotesize,
linenos
]
{python}
#!/usr/bin/python3

#EXERCICE 1
alph="abcdefghijklmnopqrstuvwxyzàçéèù"
échantillonfr="Quand, ..."
chif="IPLU'QVnÐ 'Ç'[VP3'SL'TLZZHNL'JHJOÐ'LZ['A'HGT1GXJ'("

def minuscul(texte):
	tmin=texte.lower()
	return [c for c in tmin if c in alph]

def anfreq(chaine):
	dicoccur={}
	for c in chaine:
		if c in dicoccur:
			dicoccur[c]=dicoccur[c]+1
		else:
			dicoccur[c]=1
	listocc=list(dicoccur.items())
	listocc=sorted(listocc,key=lambda t:t[1])
	listocc.reverse()
	return listocc

def replaUN(chaine,dicodéco):
	ecart=ord(dicodéco[0][1]-ord(dicodéco[0][0])
	for x in (0,len(chaine)):
		chaine[x]=chr(ord(chaine[x])+ecart)
	return chaine 
	
# Appels des fonctions (main())
statsfr=anfreq(minuscul(échantillonfr))
statsch=anfreq(chif)
llcl=[t[0] for t in statsfr]
llch=[t[0] for t in statsch]
dicodage=list(zip(llch,llcl))
#print(dicodage)
chclUN=replaUN(chif,dicodage)
\end{minted}
Cet étudiant a eu la bonne idée de faire de l'analyse fréquentielle depuis un texte écrit en français. Grâce à cette méthode il peut trouver une clef et reconstituer une partie du message. Ensuite il peut essayer de trouver la seconde clef, compte-tenu du fait qu'il connaît déjà une partie du message, en faisant le déchiffrement à la main. Cet étudiant utilise de nombreuses fonctions python élémentaires, il maîtrise la gestion de liste, les dictionnaires et les boucles.

\begin{minted}
[
frame=lines,
framesep=2mm,
baselinestretch=1.2,
bgcolor=LightGray,
fontsize=\footnotesize,
linenos
]
{python}
#!/usr/bin/python3
alphabet = 'abcdefghijklmnopqrstuvwxyz'
# La fréquence d'apparution des lettres varie bien évidemment en fonction
# de la langue et du type de texte
# Pour un texte rédigé en français, on a :
liste_frequence=['e',' ','a','s','i','t','u','r','n'] #quelques lettres 
#par ordre décroisant

chaine=input('chiffré : ')
chaine=chaine.lower()
chaine=[c for c in chaine if c in alphabet]
dico_occurences = {}

def occurence(chaine):
	for c in chaine:
	    if c in dico_occurences:
	        dico_occurences[c] = dico_occurences[c] + 1
	    else:
	        dico_occurences[c] = 1
	print(dico_occurences)

	liste_elements = list(dico_occurences.items())
	liste_elements = sorted(liste_elements,key=lambda t:t[1])
	print(liste_elements)
	dico=dict(liste_elements)
	liste_key=list(dico.keys())

	n=alphabet.index(liste_key[len(liste_key)-1])
	m=alphabet.index('e')
	if n >= m :
		decalage=n-m
	else:
		decalage=m-n
	return decalage

texte_clair=''
decalage=occurence(chaine)
decalage=26-decalage #pour déchiffrer 
alphabet_dechiffrement = alphabet[decalage:] + alphabet[:decalage]
associations = dict(zip(alphabet,alphabet_dechiffrement.upper()))
for c in chaine:
    if c in associations :
        texte_clair += associations[c]
print("############")
print("le message clair : ")
print(texte_clair.lower())
# je parviens pas à continuer !!!!!!
\end{minted}

Cet étudiant a suivi le même raisonnement que celui précédemment. Son raisonnement arrive au même résultat et en allant un peu plus loin il aurait pu arriver au résultat final. Il maitrise quand même lui aussi les listes, dictionnaires et les fonctions élémentaire de python.\\
On voit que seulement deux personnes sur quatre ont traités cet exercice, ce qui est moins que ce à quoi nous nous attendions.

\subsubsection{Second exercice}
\begin{minted}
[
frame=lines,
framesep=2mm,
baselinestretch=1.2,
bgcolor=LightGray,
fontsize=\footnotesize,
linenos
]
{python}
#!usr/bin/python3
import re
A = """155c04970e1747172e4e070e044e17776e076526
0417474e27330474074e2767072e4e070e9726044f2f776e0407e704272f
0f042e76074e970e4f27047623"""
n = len(A)

def hexad(n): #fonction qui prend transforme un nombre en haxa et 
 # le passe en entier  
	h = 0
	for i in range(len(n)):
		if n[i] == 'a':
			h += 10 * 16**(1-i)
		elif n[i] == 'b':
			h += 11 * 16**(1-i)
		elif n[i] == 'c':
			h += 12 * 16**(1-i)
		elif n[i] == 'd':
			h += 13 * 16**(1-i)
		elif n[i] == 'e':
			h += 14 * 16**(1-i)
		elif n[i] == 'f':
			h += 15 * 16**(1-i)
		else:
			h += int(n[i]) * 16**(1-i)
	return h

def bbinaire(x):
	b = str(bin(x))
	b = b[2:]
	n = len(b)
	if n<8:
		b = (8-n)*'0' + b
	return b

def echange(n):
	if n< 10000000 and n%10 == 1:
		n= n + 10000000 -1
	elif n>1000000 and n%10 == 0:
		n = n - 1000000 +1
	return n

B = []
m = len(B)

for i in range(0,n-1,2):
	B.append(bbinaire(hexad(A[i:i+2])))

#for j in range(m):
#mon but était de faire le chemin inverse avec mes fonctions 

print(B)
\end{minted}

Ici notre étudiant a créer des fonctions qui n'ont pas réellement d'intérêt puisqu'il existe des fonctions élémentaires en python qui permettent, par exemple, de passer d'un entier à sa représentation hexadécimale et vice-versa. La fonction \textsf{bbinaire} permet d'utiliser des octets. Quant à la fonction \textsf{echange} elle n'a aucun intérêt. L'étudiant a essayé quand même de faire quelque chose mais ça reste très loin du résultat final.
\begin{minted}
[
frame=lines,
framesep=2mm,
baselinestretch=1.2,
bgcolor=LightGray,
fontsize=\footnotesize,
linenos
]
{python}
#!/usr/bin/python3
import sys , re
liste1 = []
liste2 = []
message = """155C04970E1747172E4E070E044E17776E0
765260417474E27330474074E2767072E4E070E9726044F2F776e040
7e704272f0f042e76074e970e4f27047623"""

p = re.compile('\w\w')
liste1 = p.findall(message)

for lettre in liste1 :
	liste2 += bin(int(lettre, 16))

print (liste2)
\end{minted}
Cet étudiant a utilisé les expressions régulières pour récupérer deux à deux les caractères du chiffrés. Il les convertit ensuite en binaire.

\begin{minted}
[
frame=lines,
framesep=2mm,
baselinestretch=1.2,
bgcolor=LightGray,
fontsize=\footnotesize,
linenos
]
{python}
#!/usr/bin/python3
message = input("Entrez le message à envoyer? : ")
l=len(message)
liste=[]
def gest_caractere(message):
	blocs=[[message[2*i:2*i+2]] for i in range(0,l//2)] 
	#on construit des blocs de 2 elements
	print(blocs)

	for i in range(0,len(blocs)):
		liste.append(blocs[i][0])
	print(liste)
	liste_bin=[bin(int(liste[i], 16)) for i in range(0,len(liste))]
	print(liste_bin)
	#le temps ??????
gest_caractere(message)
\end{minted}
Cet étudiant a eu la bonne idée de découper le message en blocs de deux caractères. Il les stocke dans une liste qu'il convertit ensuite en une liste de nombre binaires. Les lignes 10 à 12 n'ont aucun intérêt. Il pouvait totalement travailler avec la liste \textsf{blocs}. Toutefois il n'aura effectuer que deux opérations sur la dizaine demandée.\\

L'objectif de cet exercice était pourtant clair, mais il a du susciter beaucoup d'incompréhensions vis à vis de nos étudiants qui ont utilisés beaucoup de fonctions inutiles.

\subsection{Sujet n°2}
\subsubsection{Premier exercice}
\begin{minted}
[
frame=lines,
framesep=2mm,
baselinestretch=1.2,
bgcolor=LightGray,
fontsize=\footnotesize,
linenos
]
{python}
#!usr/bin/python3
import os, sys
phrase = 'Il fait beau ce matin.'
phrase = phrase.upper()
PH = []
for u in phrase:
	PH.append(ord(u))
def premf(A):
	for i in range(len(A)):
		if A[i]%3 == 0:
			A[i] += (i+13)%i
	return A
premf(PH)
ch = ""
for i in range(len(PH)):
	ch += chr(PH[i])

CH = bytes(ch, 'utf-8')
ME = []
for i in CH:
	if i%2 ==0:
		ME.append(i^5)
	elif i%3==0:
		ME.append(i|6)
	elif i%5==0:
		ME.append(i&2)
	else:
		ME.append(i)

for i in range(len(ME)):
	ME[i] = chr(ME[i])

mes_fin = ''.join(ME)
print(mes_fin)

# j'obtiens II%CI%GI%CW%MaI[+ 
\end{minted}
Cet étudiant a bien comprit le but de cet exercice. Il utilise dans le bon ordre les différentes fonctions et a comprit de manière générale les étapes à suivre, malgré quelques erreurs. Il n'a pas utilisé les bons opérateurs logiques pour obtenir le bon chiffré.
\begin{minted}
[
frame=lines,
framesep=2mm,
baselinestretch=1.2,
bgcolor=LightGray,
fontsize=\footnotesize,
linenos
]
{python}
#!/usr/bin/python3
#Exercice 1
message = 'Il fait beau ce matin.'
msg = message.upper()
L = []
for lettre in msg:
    L.append(ord(lettre))
print(L)
L_prime = []
for nombre in L:
    if int(nombre)%3 == 0:
        L_prime.append((int(nombre)+13)%int(nombre))
    else:
        L_prime.append(nombre)
print(L_prime)
#subsidiaire
#chaine = ''
#for nombre in L_prime:
#    c = chr(nombre)
#    print(c)
#    chaine += c    
#subsidiaire
L_bin = []
print(bin(67))

for nombre in L_prime:
    b = bin(nombre)
    if int(nombre)%2 == 0:
        L_bin.append(b[2:len(b)]^101)
    elif int(nombre)%3 == 0:
        L_bin.append(b[2:len(b)] or 110)
    elif int(nombre)%5 == 0:
        L_bin.append(b[2:len(b)] and 10)
    else:
        L_bin.append(b[2:len(b)])
   
print(L_bin)
# NON FINI
\end{minted}

Cet étudiant a lui aussi bien suivit l'ordre de l'algorithme défini mais il n'a pas terminé son code. Il avait bien comprit le but de l'exercice et connaissait ses opérateurs logiques. Toutefois il a fait des erreurs comme la modification des octets suivant la parité de l'index du tableau.
\begin{minted}
[
frame=lines,
framesep=2mm,
baselinestretch=1.2,
bgcolor=LightGray,
fontsize=\footnotesize,
linenos
]
{python}
#!/usr/bin/python3
import sys
message = "Il fait beau ce matin."
message = message.upper()
tableau=[]

for i in range(0, len(message)) :
	tableau.append(ord(message[i]))
	if(tableau[i]%3==0):
		tableau[i]+= (13+i)%i
for i in range(0, len(message)-1) :
	if(i%2==0):
		tableau[i] = tableau[i]^5
	if(i%3==0):
		newl=""
		l = bin(tableau[i])
		for j in range(0, len(l)-3):
			newl += l[j]
		newl +="11"
		newl += l[len(l)-1]
		tableau[i] = int(newl, 2)
	if(i%5==0):
		newl=""
		l = bin(tableau[i])
		for j in range(0, len(l)-3):
			newl += l[j]
		newl +="0"
		newl += l[len(l)-2]
		newl +="0"
		tableau[i] = int(newl, 2)
l = ""
for i in range(0, len(message)-1) :
	l += chr(tableau[i])

print(l) #JL%FDHV BO@U'CW"HAfIZ
\end{minted}
Cet étudiant a presque le bon résultat. Sa méthodologie est bonne mais ses opérations logiques sont mal effectuées. En effet, il faut les faire avec des bits. La chaine "110" n'est pas le bit "110". Avec un peu plus de rigueur et d'attention cet étudiant aurait obtenu le bon chiffré. Il est quand même important de souligner que tout le reste est correct.
\pagebreak
\begin{minted}
[
frame=lines,
framesep=2mm,
baselinestretch=1.2,
bgcolor=LightGray,
fontsize=\footnotesize,
linenos
]
{python}
#!/usr/bin/python3

message='il fait beau ce matin'
message=message.upper()
#print(message)
tableau=[ord(i) for i in message]
print(tableau)
for i in range(0,len(tableau)):
	if tableau[i]%3==0:
		tableau[i]=tableau[i]+(13+i)%i 
#print(tableau)
tableau_caract=[chr(i) for i in tableau]
print(tableau_caract)
chaine=''.join(tableau_caract)
print(chaine)
tableau_oct=[bin(ord(i))[2:] for i in chaine]
print(tableau_oct[0])
for i in range(0,len(tableau_oct)):
	if tableau[i]%2==0:
		tableau[i]=tableau[i]^int('101',2)
	if tableau[i]%3==0:
		tableau[i]=tableau[i]|int('110',2)
	if tableau[i]%5==0:
		tableau[i]=tableau[i]&int('010',2)
#print(tableau)
chaine_f=[chr(i) for i in tableau]
chaine_f=''.join(chaine_f)
print(chaine_f)
\end{minted}

Cet étudiant a une excellente méthodologie, il a bien suivit l'énoncé. La seule erreur notable est la mauvaise écriture des opérateurs logiques en python. En dehors de cet erreur, le code est totalement correct.

\begin{minted}
[
frame=lines,
framesep=2mm,
baselinestretch=1.2,
bgcolor=LightGray,
fontsize=\footnotesize,
linenos
]
{python}
#!/usr/bin/python3
#EXERCICE 1
txt="Il fait beau ce matin."
#1
TXT=txt.upper()
#2 et 3
tab=[]

for i in range(0,len(txt)):
	tab.append(ord(txt[i]))
	if tab[i]%3==0:
		tab[i]=tab[i]+((13+i)%i)
#4
rst=""

for i in range(0,len(tab)):
	rst=rst+chr(tab[i])
#5
loct=list(map(bin,rst.encode('utf-8')))
#6
def tabmodif(ta):é
	for i in range(0,len(ta)):
		if i%2==0:
			ta[i]=ta[i]^101
		if i%3==0:
			ta[i]=ta[i] (ou inclusif) 110
		if i%5==0:
			ta[i]=ta[i] (et) 010
	return ta
#7
#Si fait avec la liste du 3)
rslt=""

for i in range(0,len(tab))
	rslt=rslt+chr(tab[i])

#Si fait avec liste du 5)
rsltat=""
for i in range(0,len(loct)):
	rsltat=rsltat+bytes([int(loct[i],2)]).decode('utf-8')
\end{minted}
Cet étudiant était sur la bonne piste mais il n'a pas été au bout de la réflexion et a manqué de rigueur sur la rédaction du code. On remarque qu'il ne connait pas les opérateurs logiques ni la conversion chaîne de caractères - bits. Même après "correction" du code, le résultat n'est pas du tout celui attendu.

\pagebreak

\subsubsection{Second exercice}
\begin{minted}
[
frame=lines,
framesep=2mm,
baselinestretch=1.2,
bgcolor=LightGray,
fontsize=\footnotesize,
linenos
]
{python}
#!usr/bin/python3

import os, sys

ph="FC?E^FK?YZ^J?XZ?B^KFM-"
S = bytes(ph,'utf-8')
A = []
for u in S:
	A.append(chr((u^21)-10))

phrase = ''.join(A)
print(phrase)

# résultat = IL FAIT BEAU CE MATIN.

#pour arriver à ce résultat j'ai stocké ordre UTF-8 
#de chaque caractère dans S. Ensuite fais les opération inverses du 
#chiffrement (le XOR avec 10101 est devenu XOR avec 21 qui est l'écriture décimal de 10101)
#j'ai reconverti tous les entiers en caractères puis j'ai concaténé la chaine
\end{minted}

Cet étudiant a réussi l'exercice. Il a même prit l'initiative d'utiliser le \textsf{xor} entre entiers, ce qui fonctionne avec python. Il obtient le bon résultat ce qui est satisfaisant.

\begin{minted}
[
frame=lines,
framesep=2mm,
baselinestretch=1.2,
bgcolor=LightGray,
fontsize=\footnotesize,
linenos
]
{python}
#!/usr/bin/python3
#Exercice 2
message = 'FC?EˆFK?YZˆJ?XZ?BˆKFM-'
L = []
for lettre in message:
    L.append(ord(lettre) + 10)
print(L)
msg = ''
for i in L:
    msg += chr(i)
print(msg)
# NON FINI
\end{minted}
Cet étudiant n'a effectué que les trois premières étapes, sûrement par faute de temps. Il était toutefois sur la bonne voie. Les premières lignes de codes sont justes.

\pagebreak

\begin{minted}
[
frame=lines,
framesep=2mm,
baselinestretch=1.2,
bgcolor=LightGray,
fontsize=\footnotesize,
linenos
]
{python}
#!/usr/bin/python3
import sys
message = "FC?E^FK?YZ^J?XZ?B^KFM"
dechiffre = ""
for i in range(0, len(message)) :
	dechiffre += chr((ord(message[i])^21)-10)
print(dechiffre) #IL FAIT BEAU CE MATIN
\end{minted}
Cet étudiant a également réussi l'exercice et de façon encore plus compacte. Il n'y a pas grand chose à redire, sa maîtrise du \textsf{xor} entre entiers est intéressante.\\
\begin{minted}
[
frame=lines,
framesep=2mm,
baselinestretch=1.2,
bgcolor=LightGray,
fontsize=\footnotesize,
linenos
]
{python}
#Chiffrement:
chif='FC?E^FK?YZ^J?XZ?B^KFM-'
tabl=[]
chaine=""
for i in range(0,len(chif)):
	tabl.append(ord(chif[i])+10)
	chaine=chaine+chr(tabl[i])
lisoct=list(map(bin,chaine.encode('utf-8')))
for i in range(0,len(lisoct)):
	lisoct[i]=lisoct[i]^10101
retout=""
for i in range(0,len(lisoct)):
	retour=retour+bytes([int(lisoct[i],2)]).decode('utf-8')
#Déchiffrement
byt=list(map(bin,chif.encode('utf-8')))
inter=""
clair=""
for i in range(0,len(byt)):
	byt[i]=byt[i]^10101
	inter=inter+bytes([int(byt[i],2)]).decode('utf-8')
t=[]
for i in range(0,len(inter)):
	t.append(ord(inter[i]))
	t[i]=t[i]-10
	clair=clair+chr(t[i])
\end{minted}

Cet étudiant a essayé de coder la méthode de chiffrement, ce qui était inutile. Peut-être espérait-il que ça l'aiderait. Ceci étant, sa méthode de chiffrement ne marche pas. En effet, il a tenté d'effectuer un \textsf{xor} entre deux types différents à la ligne $10$.\\
Pour le déchiffrement il a effectué la même erreur, c'est à dire un \textsf{xor} entre deux types différents (à la ligne $19$) donc sa méthode ne marche pas. Cela dit, hormis cette erreur importante, le reste du programme semble correct.\\
On voit tout de même que c'était un exercice relativement rapide à faire. Seules 7 lignes de code suffisaient et il était probablement possible de faire légèrement plus court.
\subsubsection{Troisième exercice}
A notre plus grand regret cet exercice n'a pas été traité. Cela aurait été un bon moyen de départager les étudiants et c'est celui qui ressemblait le plus à un CTF (Capture The Flag).
\subsection{Conclusion}
\end{document}